\section*{अध्याय \ref{ch:b2:solid-liquid-diagrams} --- द्वि-घटकी ठोस--द्रव आरेख}

\begin{solution}{exo:b2:solid-liquid-diagrams:1}
\qty{1300}{\degreeCelsius} पर \omterm{def:b2:solid-liquid-diagrams:liquidus}{लिक्विडस} 0.541 पर और \omterm{def:b2:solid-liquid-diagrams:liquidus}{सॉलिडस} 0.609 पर है।
0.40: केवल द्रव। 0.58: द्रव (0.541) और \omterm{def:b2:solid-liquid-diagrams:solid-solution}{ठोस विलयन} (0.609)। 0.70:
केवल \omterm{def:b2:solid-liquid-diagrams:solid-solution}{ठोस विलयन}।
\end{solution}

\begin{solution}{exo:b2:solid-liquid-diagrams:2}
ठोस अंश $(0.58 - 0.541)/(0.609 - 0.541) = 0.57$: ठोस के \qty{1.15}{mol},
द्रव के \qty{0.85}{mol}। निकेल: ठोस में $1.15 \times 0.609 = \qty{0.70}{mol}$,
द्रव में $0.85 \times 0.541 = \qty{0.46}{mol}$; कुल \qty{1.16}{mol}
$= 2.00 \times 0.58$।
\end{solution}

\begin{solution}{exo:b2:solid-liquid-diagrams:3}
नियत दाब पर $v = n + 1 - \varphi$ (कोई अभिक्रिया नहीं)। (a) $1 + 1 - 2 = 0$।
(b) $2 + 1 - 2 = 1$। (c) द्रव और दो \omterm{def:b2:solid-liquid-diagrams:solid-solution}{ठोस विलयन}: $2 + 1 - 3 = 0$।
\end{solution}

\begin{solution}{exo:b2:solid-liquid-diagrams:4}
\qty{232}{\degreeCelsius} तक सुचारु गिरावट, टिन के जमने तक एक क्षैतिज पठार,
फिर दोबारा सुचारु गिरावट, कमरे के ताप के निकट धीमी। पठार पर
नियत दाब पर द्रव और ठोस सह-अस्तित्व में हैं, \omterm{def:b2:equilibrium-shifts:variance}{स्वातंत्र्य कोटि} 0: हटाई गई ऊष्मा
क्रिस्टलन से आती है, और अंतिम द्रव के जमने तक ताप नहीं बदल
सकता।
\end{solution}

\begin{solution}{exo:b2:solid-liquid-diagrams:5}
बेन्ज़ीन: $\ln x_B = -(9870/8.314)(1/269.6 - 1/278.64) = -0.1429$, $x_B =
0.867$। नैफ़्थलीन: $\ln x_N = -(19\,100/8.314)(1/269.6 - 1/353.2) = -2.017$,
$x_N = 0.133$। योग 1.000 है: \qty{269.6}{K} ($\qty{-3.5}{\degreeCelsius}$)
गलनक्रांतिक ताप है, और गलनक्रांतिक द्रव में 0.133 नैफ़्थलीन है।
\end{solution}

\begin{solution}{exo:b2:solid-liquid-diagrams:6}
द्रव सीसे वाले पक्ष के \omterm{def:b2:solid-liquid-diagrams:liquidus}{लिक्विडस} तक ठंडा होता है, जहाँ (Pb) के क्रिस्टल प्रकट होते हैं;
द्रव टिन में समृद्ध होता जाता है, \qty{182.2}{\degreeCelsius}
और 62.1~\% पर गलनक्रांतिक तक। वहाँ शेष द्रव \omterm{def:b2:solid-liquid-diagrams:eutectic}{गलनक्रांतिक मिश्रण} में जम जाता है।
गलनक्रांतिक का द्रव्यमान अंश: $(30 - 18.91)/(62.13 - 18.91) = 0.26$। और
ठंडा करने पर (Pb) में टिन की विलेयता घटती है (खंडित रेखा): (Pb) क्रिस्टलों के भीतर
थोड़ा (Sn) अवक्षेपित होता है।
\end{solution}

\begin{solution}{exo:b2:solid-liquid-diagrams:7}
द्रव्यमान के अनुसार 23.3~\% NaCl: प्रति किलोग्राम जल $23.3/76.7 = \qty{0.304}{kg}$
लवण। \qty{-25}{\degreeCelsius} पर, गलनक्रांतिक से नीचे, अनुपात चाहे जो हों, कोई
द्रव अस्तित्व में नहीं रह सकता: नमक और बर्फ़ ठोसों के रूप में अगल-बगल रहते हैं।
\end{solution}

\begin{solution}{exo:b2:solid-liquid-diagrams:8}
$x_{\mathrm B} = 0.5$, E$_1$ और यौगिक के बीच है: $\mathrm{AB_2}$
पहले क्रिस्टलित होता है, द्रव E$_1$ की ओर बढ़ता है, वहाँ A और $\mathrm{AB_2}$ के
गलनक्रांतिक में जमता है; अंत में A + $\mathrm{AB_2}$।
$x_{\mathrm B} = 0.9$, E$_2$ और B के बीच है: B पहले क्रिस्टलित होता है, द्रव
E$_2$ तक पहुँचता है; अंत में $\mathrm{AB_2}$ + B।
\end{solution}

\begin{solution}{exo:b2:solid-liquid-diagrams:9}
मंडल के पिछले किनारे पर द्रव ऐसे ठोस में जमता है जिसमें उसके ताँबे के मोल अंश का केवल 0.78
गुना होता है: अस्वीकृत ताँबा द्रव में रहता है, जो आगे बढ़ता है और
समृद्ध होता जाता है, इसलिए ताँबा छड़ के सिरे की ओर बुहार दिया जाता है।
1 के निकट अनुपात के साथ हर पारण अशुद्धि का केवल एक अंश हटाता है:
कई पारण चाहिए, हर एक पिछले पारण से बची छड़ से आरंभ करता है।
\end{solution}

\begin{solution}{exo:b2:solid-liquid-diagrams:10}\emergencystretch=3em
व्युत्पत्ति अध्याय की तरह। तनु द्रव के लिए $\ln x_A \approx -x_B$, $1/T -
1/T^* \approx -\Delta T/T^{*2}$, अतः $x_B = \Delta_{\text{गलन}}H_A\Delta T/(RT^{*2})$;
$x_B \approx n_B/n_A = bM_A$ के साथ $\Delta T = (RT^{*2}M_A/\Delta_{\text{गलन}}H_A)\,b$।
बेन्ज़ीन: $K_f = 8.314 \times 278.64^2 \times 0.07811/9870 = \qty{5.11}{K.kg/mol}$।
\end{solution}

\begin{solution}{exo:b2:solid-liquid-diagrams:11}
प्रति \qty{100}{g}: Mg के $19.0/24.31 = \qty{0.782}{mol}$, Pb के $81.0/207.2 =
\qty{0.391}{mol}$; अनुपात $2.00$: $\ce{Mg2Pb}$।
\end{solution}

\begin{solution}{exo:b2:solid-liquid-diagrams:12}
पठार गलनक्रांतिक द्रव के द्रव्यमान के समानुपाती है: अज्ञात में उसका
$120/300 = 0.40$ है। सीसे वाले पक्ष पर $w = 18.91 + 0.40 \times (62.13 -
18.91) = 36.2~\%$ टिन; टिन वाले पक्ष पर $w = 96.59 - 0.40 \times (96.59 - 62.13) =
82.8~\%$। \omterm{def:b2:solid-liquid-diagrams:cooling-curve}{शीतलन वक्र} का पहला मोड़ (36~\% वाली मिश्रधातु के लिए ऊँचा,
जिसका \omterm{def:b2:solid-liquid-diagrams:liquidus}{लिक्विडस} खड़े सीसा-पक्ष पर है) या एक सूक्ष्मचित्र (प्राथमिक (Pb) या
प्राथमिक (Sn) क्रिस्टल) उन्हें अलग पहचानता है।
\end{solution}

\begin{solution}{pb:b2:solid-liquid-diagrams:1}
\textbf{1.} सीसा \qty{327.5}{\degreeCelsius}, टिन \qty{231.9}{\degreeCelsius}।
\textbf{2.} प्रति \qty{100}{g}: टिन के $62.13/118.71 = \qty{0.5234}{mol}$ और
सीसे के $37.87/207.2 = \qty{0.1828}{mol}$; $x_{\mathrm{Sn}} = 0.5234/0.7062 =
0.741$।
\textbf{3.} बिंदु (0, 327.5), (100, 231.9), (62.13, 182.2), (18.91, 182.2),
(96.59, 182.2)।
\textbf{4.} चित्र की तरह: द्रव; L + (Pb); L + (Sn); (Pb); (Sn); (Pb) + (Sn)।
\textbf{5.} L (62.1~\% Sn) $\longrightarrow$ (Pb) (18.9~\%) + (Sn) (96.6~\%)।
\textbf{6.} $v = 2 + 1 - 3 = 0$: ताप और तीनों संघटन
नियत हैं।
\textbf{7.} सीसा-समृद्ध \omterm{def:b2:solid-liquid-diagrams:solid-solution}{ठोस विलयन} (Pb): सीसा टिन घोलता है, इसलिए द्रव के साथ
साम्य वाले ठोस में पहले से टिन होता है, उस अनुपात में जो
\omterm{def:b2:solid-liquid-diagrams:liquidus}{सॉलिडस} पर पढ़ा जाता है।
\textbf{8.} वह \omterm{def:b2:solid-liquid-diagrams:liquidus}{लिक्विडस} के अनुदिश गलनक्रांतिक की ओर बढ़ता है, टिन में समृद्ध होता हुआ,
62.13~\% तक।
\textbf{9.} 62.13~\% का द्रव और 18.91~\% टिन वाला (Pb)।
\textbf{10.} (Pb): $(62.13 - 40)/(62.13 - 18.91) = 0.512$; द्रव 0.488।
\textbf{11.} 18.91~\% का (Pb) और 96.59~\% का (Sn); (Sn): $(40 - 18.91)/(96.59 -
18.91) = 0.271$; (Pb): 0.729।
\textbf{12.} गलनक्रांतिक घटक 48.8~\% (ठीक ऊपर उपस्थित द्रव);
प्राथमिक (Pb) 51.2~\%। प्रश्न 11 का (Pb) प्राथमिक क्रिस्टलों और
गलनक्रांतिक की (Pb) पटलिकाओं का योग है।
\textbf{13.} \omterm{def:b2:solid-liquid-diagrams:liquidus}{लिक्विडस} पर मोड़, धीमी गिरावट, फिर
\qty{182.2}{\degreeCelsius} पर पठार जो गलनक्रांतिक मिश्रधातु के पठार का 0.488 गुना चलता है।
\textbf{14.} $\ln x_{\mathrm{Sn}} = -(7030/8.314)(1/T - 1/505.1)$।
\textbf{15.} $\ln x_{\mathrm{Bi}} = -(11\,300/8.314)(1/T - 1/544.6)$।
\textbf{16.} $x_{\mathrm{Sn}}(T) + x_{\mathrm{Bi}}(T) = 1$।
\textbf{17.} \qty{110}{\degreeCelsius} पर: $0.587 + 0.349 = 0.936$;
\qty{120}{\degreeCelsius} पर: $0.621 + 0.382 = 1.003$;
\qty{130}{\degreeCelsius} पर: $0.655 + 0.417 = 1.072$। योग लगभग
\qty{119.5}{\degreeCelsius} पर 1 तक पहुँचता है: निकटतम डिग्री तक \qty{120}{\degreeCelsius}।
\textbf{18.} $x_{\mathrm{Bi}} = 0.38$; द्रव्यमान में $0.38 \times 208.98/(0.38 \times
208.98 + 0.62 \times 118.71) = 0.52$।
\textbf{19.} मॉडल \qty{120}{\degreeCelsius} और 52~\% बिस्मथ देता है,
मूल्यांकित आरेख \qty{138.8}{\degreeCelsius} और 57~\%: लगभग
\qty{19}{\degreeCelsius} बहुत नीचा।
\textbf{20.} क्रिस्टलित होने वाला टिन शुद्ध नहीं, बल्कि बिस्मथ वाला एक
विलयन है: उसका \omterm{def:b2:chemical-potential:chemical-potential}{रासायनिक विभव} शुद्ध टिन से कम है, ठोस
अधिक स्थायी है, और दिए गए द्रव संघटन के लिए वह द्रव से उच्चतर ताप पर क्रिस्टलित
होता है। \omterm{def:b2:solid-liquid-diagrams:liquidus}{लिक्विडस} की टिन शाखा ऊपर उठती है, और वह
बिस्मथ शाखा से ऊँचे पर मिलती है (अनादर्श द्रव अपना स्वयं का संशोधन जोड़ता है)।
\textbf{21.} वह एक ही ताप पर, निकाय के निम्नतम ताप पर, पिघलता और जमता है:
कोई लेईदार परिसर नहीं जिसमें हिलाया गया जोड़ चटक जाए।
\textbf{22.} सीसा विषैला है, और इलेक्ट्रॉनिक अपशिष्ट पर्यावरण में पहुँचता है।
\textbf{23.} ऊष्मा-संवेदी घटकों को निम्नतर ताप पर सोल्डर किया जा सकता है;
परंतु उपयोग में जोड़ निम्नतर ताप पर नरम पड़ता है, और बिस्मथ उसे
भंगुर बनाता है।
\textbf{24.} \qty{570}{g} बिस्मथ और \qty{430}{g} टिन।
\textbf{25.} आदर्श-द्रव मॉडल $\boldsymbol{T_E
\approx \qty{120}{\degreeCelsius}}$ पर गलनक्रांतिक का पूर्वानुमान देता है, जबकि
मूल्यांकित मान \qty{138.8}{\degreeCelsius} है।
\end{solution}
